\documentclass[a4paper,10pt]{report}\input{../0.Préambule/Préambule_cours_prof}\begin{document}

\newpage
\pagestyle{empty}
\section*{La démonstration en mathématiques}

\noindent Pour rédiger une démonstration : 
\begin{enumerate}
\item[•] On justifie rapidement les déductions simples

\medskip
\item[•] Pour les déductions importantes, \\on utilise le plan :
\end{enumerate} 

\vspace{-1.4cm}
\begin{flushright}
\boxed{\begin{minipage}{10cm}
\noindent \begin{enumerate}
\item[•] \textbf{Je sais que} : \textit{noter les hypothèses nécessaires à la propriété}
\item[•] \textbf{J'utilise la propriété }: \textit{....... citer la propriété}
\item[•]\textbf{Je conclus que} : \textit{........ conclure}
\end{enumerate}
\end{minipage}}
\end{flushright}


\vspace{-7mm}
\setcounter{cex}{0}
\begin{ex}
$\mathcal{C}$ et $\mathcal{C}’$ sont deux cercles de centre $O$.

\noindent $[DM]$ est un diamètre du cercle $\mathcal{C}$ et $[EI]$ est un diamètre du cercle $\mathcal{C}’$.
\begin{enumerate}
\item Faire une figure.
\item Démontrer que $DEMI$ est un parallélogramme.
\end{enumerate}
\end{ex}

\vspace{1cm}
\section*{La démonstration en mathématiques}

\noindent Pour rédiger une démonstration : 
\begin{enumerate}
\item[•] On justifie rapidement les déductions simples

\medskip
\item[•] Pour les déductions importantes, \\on utilise le plan :
\end{enumerate} 

\vspace{-1.4cm}
\begin{flushright}
\boxed{\begin{minipage}{10cm}
\noindent \begin{enumerate}
\item[•] \textbf{Je sais que} : \textit{noter les hypothèses nécessaires à la propriété}
\item[•] \textbf{J'utilise la propriété }: \textit{....... citer la propriété}
\item[•]\textbf{Je conclus que} : \textit{........ conclure}
\end{enumerate}
\end{minipage}}
\end{flushright}


\vspace{-7mm}
\setcounter{cex}{0}
\begin{ex}
$\mathcal{C}$ et $\mathcal{C}’$ sont deux cercles de centre $O$.

\noindent $[DM]$ est un diamètre du cercle $\mathcal{C}$ et $[EI]$ est un diamètre du cercle $\mathcal{C}’$.
\begin{enumerate}
\item Faire une figure.
\item Démontrer que $DEMI$ est un parallélogramme.
\end{enumerate}
\end{ex}

\vspace{1cm}
\section*{La démonstration en mathématiques}

\noindent Pour rédiger une démonstration : 
\begin{enumerate}
\item[•] On justifie rapidement les déductions simples

\medskip
\item[•] Pour les déductions importantes, \\on utilise le plan :
\end{enumerate} 

\vspace{-1.4cm}
\begin{flushright}
\boxed{\begin{minipage}{10cm}
\noindent \begin{enumerate}
\item[•] \textbf{Je sais que} : \textit{noter les hypothèses nécessaires à la propriété}
\item[•] \textbf{J'utilise la propriété }: \textit{....... citer la propriété}
\item[•]\textbf{Je conclus que} : \textit{........ conclure}
\end{enumerate}
\end{minipage}}
\end{flushright}


\vspace{-7mm}
\setcounter{cex}{0}
\begin{ex}
$\mathcal{C}$ et $\mathcal{C}’$ sont deux cercles de centre $O$.

\noindent $[DM]$ est un diamètre du cercle $\mathcal{C}$ et $[EI]$ est un diamètre du cercle $\mathcal{C}’$.
\begin{enumerate}
\item Faire une figure.
\item Démontrer que $DEMI$ est un parallélogramme.
\end{enumerate}
\end{ex}
\end{document}